\section{递归地址化：派生概念由旧层读出序列生成}\label{sec:recursive-addressing}

本节仅处理一个结构性问题：在给定一族可测概念之后，如何由其有限时间窗读出递归生成更细的概念族，以及该过程是否会在可测结构层面引入“外加对象”。本文在可测动力系统的标准框架下工作（第 \ref{sec:preliminaries} 节），并证明递归地址化严格等价于对既有柱事件族取 \(\sigma\)-代数细化（命题 \ref{prop:recursive-addressing-refinement}）。从拓扑直觉看，柱事件前缀亦对应一条前缀超度量过滤：长度前缀等价于更小半径的超度量球，见第 \ref{sec:spg} 节段落 \ref{par:spg-prefix-ultrametric}。

\subsection{地址的复合生成}

设在层级 $L$ 已得到一族概念 $\mathcal{C}^{(L)}=\{C^{(L)}_1,\dots,C^{(L)}_{r_L}\}$，其中每个 $C^{(L)}_i\in\mathcal{B}$ 可测。定义派生的二元向量过程
\[
\mathbf{b}^{(L)}_t(x):=\bigl(\ind_{C^{(L)}_1}(T^t x),\dots,\ind_{C^{(L)}_{r_L}}(T^t x)\bigr)\in\{0,1\}^{r_L}.
\]
对任意长度 $\ell$ 与任意地址
\[
a=\bigl(a_0,\dots,a_{\ell-1}\bigr)\in\bigl(\{0,1\}^{r_L}\bigr)^\ell,
\]
定义其对应的派生柱事件
\[
C^{(L+1)}_{a,t}:=\{x\in X:\mathbf{b}^{(L)}_t(x)=a_0,\dots,\mathbf{b}^{(L)}_{t+\ell-1}(x)=a_{\ell-1}\}\in\mathcal{B}.
\]
令 $\mathcal{C}^{(L+1)}$ 为选定的一族此类柱事件（例如取一组有限地址集合 $A^{(L+1)}$ 对应的 $C^{(L+1)}_{a,0}$）。由此得到闭合递归：\textbf{派生概念是旧层读出序列的柱事件}，从而地址仅来自可观测数据而非外加原语。

\begin{proposition}[$\sigma$-代数细化：递归地址化不引入外加结构]\label{prop:recursive-addressing-refinement}
令 $\mathcal{G}^{(L)}$ 为由旧层概念在有限时间窗内的平移生成的 $\sigma$-代数：
\[
\mathcal{G}^{(L)}:=\sigma\bigl(T^{-t}C_i^{(L)}:\ 1\le i\le r_L,\ t\in\ZZ\bigr).
\]
则对任意地址 $a$ 与任意 $t$，派生柱事件 $C^{(L+1)}_{a,t}$ 属于 $\mathcal{G}^{(L)}$。因此递归地址化仅执行“对既有可观测事件作柱集闭包/细化”，而不引入外加可测集。
\end{proposition}

\begin{proof}
对任意固定的 \(i\) 与 \(t\)，函数 \(x\mapsto \ind_{C_i^{(L)}}(T^t x)=\ind_{T^{-t}C_i^{(L)}}(x)\) 是 \(\mathcal{G}^{(L)}\)-可测的。
因此对任意给定地址 \(a=(a_0,\dots,a_{\ell-1})\)，事件
\[
\{x:\mathbf b^{(L)}_t(x)=a_0,\dots,\mathbf b^{(L)}_{t+\ell-1}(x)=a_{\ell-1}\}
\]
可写为有限个集合 \(T^{-s}C_i^{(L)}\) 及其补集的有限交（在 \(s=t,\dots,t+\ell-1\) 上并对 \(i=1,\dots,r_L\) 展开），故属于由这些集合生成的 \(\sigma\)-代数，即属于 \(\mathcal{G}^{(L)}\)。证毕。
\end{proof}

\subsection{空值到投影的递归语义}

对任意候选地址 $a\in A^{(L+1)}$，在定义派生概念 $C_a^{(L+1)}$ 之前不进行赋值；一旦地址 $a$ 被选定，概念值即为二元投影 $\ind_{C_a^{(L+1)}}(x)$。从而层级递归满足：
\begin{itemize}
  \item “无地址 $\Rightarrow$ 空值（不定义）”；
  \item “有地址 $\Rightarrow$ 二元投影（$0/1$）”；
  \item “地址 $\Rightarrow$ 由旧层读出生成”。
\end{itemize}
本节所述空值属于\emph{语义空值}（地址缺失，读出不定义）；与之正交的\emph{协议空值}（类型系统锁定 unitary slice）与\emph{碰撞空值}（侧信息预算不足导致不可唯一重建）见 \S\ref{subsec:xi-pw-type-safety-null} 及命题 \ref{prop:xi-null-three-way}。在更统一的口径下，可把“完备地址”打包为多轴前缀（定义 \ref{def:multi_axis_address}），其中碰撞空值对应残差轴不足，其最小侧信息预算由定理 \ref{thm:address-defect-law} 给出。
该语义可与布尔代数闭包兼容，也可在算子代数层面以投影算子及其乘积实现（见附录 \ref{app:op-algebra}）。

\subsection{前缀站点与局部截面：\texorpdfstring{$\mathsf{NULL}$}{NULL} 的内禀刻画}\label{subsec:prefix-site-null}
本小节给出一个与“地址先于取值”严格兼容的形式化：把地址空间视为由前缀偏序生成的站点（cylinder site），并把“协议可接受且单值可核验的读出/证书”组织为其上的（预）层对象。其用途不在于引入新的测度论结构，而在于把空值（\(\mathsf{NULL}\)）提升为一个清晰的存在性障碍：\emph{在给定地址处不存在协议允许的局部截面}，从而避免把 “不定义” 误读为数值 \(0\) 的歧义。

\begin{definition}[前缀站点与柱覆盖]\label{def:prefix-site}
为简明起见仅以二元字母表 \(\Sigma=\{0,1\}\) 书写（对一般有限字母表 \(\Sigma\) 完全同型）。令 \(\mathcal{A}:=\Sigma^{<\infty}\) 为全体有限地址，并以前缀关系 \(a\preceq b\) 赋偏序。
令 \(\mathcal{A}_{\mathrm{adm}}\subseteq\mathcal{A}\) 为协议所允许的地址集合（在本节的递归构造中，其原型可取为各层选定的有限地址族 \(A^{(L)}\) 及其在 refinement 下的闭包）。
对任意 \(a\in\mathcal{A}_{\mathrm{adm}}\)，令柱事件 \(C_a\subseteq X\) 如第 \ref{sec:spg} 节定义（亦即 \(C_a=\omega_\infty^{-1}([a])\)，其中 \([a]\subseteq\Omega_\infty\) 为符号柱集）。
\par
称有限族 \(\{b_i\succeq a\}_{i=1}^n\subseteq\mathcal{A}_{\mathrm{adm}}\) 为 \(a\) 的一个\textbf{柱覆盖}，若 \(C_a=\bigcup_{i=1}^{n} C_{b_i}\)。
该覆盖关系与段落 \ref{par:spg-prefix-ultrametric} 的前缀拓扑一致：\(\Omega_\infty=\varprojlim \Omega_k\) 的 clopen 基由柱集 \([a]\) 给出，因而“前缀细化”与“覆盖分解”为同一对象的两种等价描述。
\end{definition}

\begin{definition}[协议可接受读出的（预）层]\label{def:prefix-sheaf-readout}
固定一套“可见域 + 检查器”的读出协议，其抽象口径见定义 \ref{def:xi-visible-read-null}（在此仅使用其“读出可为部分映射、并以 \(\mathsf{NULL}\) 表示不定义”的结构）。
对每个 \(a\in\mathcal{A}_{\mathrm{adm}}\) 定义集合
\[
\mathscr{F}(a):=\{\text{在柱事件 }C_a\text{ 上协议允许、且在类型意义下为单值可核验的读出/证书对象}\},
\]
并对 refinement \(b\succeq a\) 定义限制映射
\(\rho_{a\leftarrow b}:\mathscr{F}(a)\to\mathscr{F}(b)\)
为把读出/证书限制到子柱 \(C_b\)（等价地：在事件层面对更深前缀条件化）。
\end{definition}

\begin{proposition}[局部截面存在性与 \texorpdfstring{$\mathsf{NULL}$}{NULL}]\label{prop:null-as-local-section-obstruction}
在定义 \ref{def:prefix-site}--\ref{def:prefix-sheaf-readout} 的口径下，“地址先于取值”的语义可等价表述为：
\[
\text{地址 }a\text{ 的读出在类型层非 }\mathsf{NULL}
\quad\Longleftrightarrow\quad
\mathscr{F}(a)\neq\varnothing.
\]
并且在递归地址化的闭包口径下（命题 \ref{prop:recursive-addressing-refinement}），\(\mathscr{F}\) 对任意有限柱覆盖满足粘接唯一性：若 \(\{b_i\succeq a\}\) 覆盖 \(a\)，且一族局部对象 \(s_i\in\mathscr{F}(b_i)\) 在重叠处相容（即在任意交 \(C_{b_i}\cap C_{b_j}\) 上的限制一致），则存在唯一 \(s\in\mathscr{F}(a)\) 使得 \(s|_{C_{b_i}}=s_i\)。
\end{proposition}

\begin{proof}
第一条等价仅是将 “不定义（\(\mathsf{NULL}\)）” 从数值层改写为对象层的存在性断言：读出/证书之所以可核验，要求其在给定地址处确有协议允许的局部对象可供引用；当且仅当该局部对象存在时，读出在类型层被视为非 \(\mathsf{NULL}\)。
\par
对粘接唯一性：柱覆盖是有限并集分解，且柱事件在 refinement 下封闭（由定义）；递归地址化又保证所有派生对象仍落在既有可观测柱事件所生成的 \(\sigma\)-代数内（命题 \ref{prop:recursive-addressing-refinement}）。
因此若局部对象在重叠处一致，则它们在集合层的拼接由并集给出且唯一；协议可接受性在限制映射下保持，故粘接得到的对象仍属于 \(\mathscr{F}(a)\) 并满足唯一性。
\end{proof}

\begin{corollary}[\texorpdfstring{$\mathsf{NULL}$}{NULL} 的三分解（局部截面障碍版）]\label{cor:null-trichotomy-local-section}
在上述站点/（预）层表述下，\(\mathsf{NULL}\) 的三类来源（定义 \ref{def:xi-visible-read-null} 与命题 \ref{prop:xi-null-three-way}）可被统一为三类互不混淆的局部截面障碍：
\begin{enumerate}
  \item \textbf{语义空值：} \(a\notin\mathcal{A}_{\mathrm{adm}}\)，即协议未提供该地址对象；此时在站点上无对象可供取截面，空值由“缺乏指称”直接给出。
  \item \textbf{协议空值：} \(a\in\mathcal{A}_{\mathrm{adm}}\) 但可见域/类型系统拒绝在 \(C_a\) 上建立读出（典型地：可见域被锁定为某一切片而拒绝额外坐标轴），从而 \(\mathscr{F}(a)=\varnothing\)。
  \item \textbf{碰撞空值：} \(a\in\mathcal{A}_{\mathrm{adm}}\) 且通过检查器，但在“单值可核验读出”的类型约定下，若底层读出存在多对一压缩而未外置使扩展读出单射化的残差轴，则在该类型中仍等价表现为 \(\mathscr{F}(a)=\varnothing\)；其最小侧信息预算由命题 \ref{prop:xi-null-three-way} 与定理 \ref{thm:address-defect-law} 给出。
\end{enumerate}
因此，“\(\mathsf{NULL}\)”并非字符串的数值性质，而是站点对象与（预）层纤维在协议意义下的结构性空缺：合法性由 \(\mathscr{F}(a)\) 的非空性判定。
\end{corollary}
